\section{Discussion}
\label{sec:discussion}

\paragraph{What the evidence supports.} Three statements are well supported. First, Google's benchmark numbers are real: they
reproduce with the public code and weights. Second, TimesFM~3's advantage over the pre-specified classical ensemble carries over to
data that did not exist when the model was released, with a margin (16\% in MASE, 29\% in SQL on the contrast without covariates)
that is large in practical terms and somewhat smaller, for the median forecast, than on the benchmark. The SQL margin depends on how windows
are weighted and on the river baseline (Section~\ref{sec:robustness}): with every window weighted equally it is 20\%, and with a
re-specified ensemble for river flow 22\%. Third, the advantage over the strongest other
foundation models is small: about 5\% over Chronos-2 on post-release data, and over TimesFM~2.5 resolved only in SQL. On the
benchmark the published lead over Chronos-2 holds; it comes from the tasks where cross-learning matters, while on the other tasks
the two are at parity. Published comparisons among foundation models should state the inference options (cross-learning, symmetric
averaging) each model was run with.

\paragraph{Exposure and memorisation.} We find no resolved difference between pages and windows inside the documented pageview corpus
and the same pages three years later; the estimate (\Htwo) points in the predicted direction and the interval is compatible with an
advantage of up to about \HtwoMde\%. The probe is narrow (one source, one model family, a corpus
described only by its sources), and it cannot exclude smaller effects or memorisation of series types rather than windows. Because the 2026 windows are harder than the 2023 windows and the difference in differences varies with page difficulty
(Section~\ref{sec:robustness}), the probe does not identify an effect of memorisation; it is not evidence that exposure does not help.

\paragraph{When to use it, and when not to (RQ7).} The controlled data separate the situations sharply.
\begin{itemize}
\item \emph{Short histories, long seasonality, several seasonal cycles, new series}: the foundation models are far ahead of
classical methods, which lack the information to estimate seasonality. This is the main practical case for zero-shot forecasting.
\item \emph{Strongly structured, long series} (load, prices, weather): in our domains and weeks, a 10--20\% gain over a classical
ensemble at the cost of a GPU; the gain over the long context is real but modest (4\% from 2{,}016 to about 15{,}000 hourly points).
\item \emph{Near-random-walk series} (exchange rates, interest rates, commodity prices): in our data its MASE was 11\% higher than that of the naive forecast, as were those of the
other foundation models; use a random walk as the benchmark.
\item \emph{Recent regime changes and saturating growth}: check the forecast against a simple baseline; these are the cases where
TimesFM~3 fails by large factors in our synthetic tests.
\item \emph{Covariates}: useful when the covariate drives the target strongly and its future is known; with real covariates of
moderate strength, the gain is small even when their observed future values are supplied.
\end{itemize}
On compute (Tier B, 4{,}136 forecast tasks in univariate and group mode, run records), TimesFM~3 in univariate mode took about 17 minutes on a free T4 GPU,
Chronos-2 about 5, TimesFM~2.5 and TiRex-2 about 20, and the benchmark flags of TimesFM~3 about 35; the three MSTL members of the
hourly classical ensemble took about an hour on eight CPU cores. Zero-shot forecasting is therefore not expensive at this scale;
the costs that matter in practice are the GPU dependency and the licence (TimesFM~3's weights are licensed for non-commercial
use only). Latency at batch sizes 1, 32 and 256 is reported in Appendix~\ref{app:hardware}.

\paragraph{Practical implications for digital marketing and demand planning.} We did not evaluate marketing or sales data. The
following are hypotheses that our results suggest for common use cases, to be validated on the user's own data with a loss that
reflects the decision. (i) \emph{Demand and inventory}: for intermittent demand the foundation models' quantiles were more
accurate than the classical methods we ran, their point forecasts were not better on squared error, and Croston-type methods did not
improve on the historic mean; short histories, where classical methods lack data to estimate seasonality, were the clearest case for zero-shot
forecasting in our synthetic cells. (ii) \emph{Sales with promotions or advertising}: in our data, observed future temperature
for load and precipitation for rivers gave no resolved gain, and the interval is compatible with gains and losses; whether a
planned promotion helps should be tested rather than assumed. (iii) \emph{Attention and traffic} (pageviews, search interest):
the models' 80\% intervals covered 72--76\% of outcomes on Wikipedia pageviews, so alert thresholds built on them may fire more
often than intended; in an exploratory analysis their upper quantile rarely reached the level of a spike before it started (9\%
of such windows for TimesFM~3), which is not a test of spike detection. (iv) \emph{Prices and financial series}: on the exchange
rates, yields and oil prices we evaluated, a random walk was the better benchmark. (v) \emph{Product launches and growth curves}:
the spurious decline on a saturating synthetic process and the failures after recent synthetic level shifts suggest running a naive
or damped-trend baseline alongside. The non-commercial licence of TimesFM~3 matters for any business use.
